\section{Πρόβλημα 4}

\subsection{Ερώτηση (i)}

Θεωρούμε το παρακάτω παίγνιο μηδενικού αθροίσματος:
$$
    \begin{bmatrix}
        a & 10 \\
        8 & d
    \end{bmatrix}
$$
όπου $0 \leq a < 10$ και $d > 10$. Για να βρούμε αν υπάρχει σημείο ισορροπίας
θα πρέπει να ελέγξουμε τα μέγιστα και ελάχιστα κάθε γραμμής και στήλης. Για τον
παίκτη των γραμμών έχουμε: $i)$ εφόσον $a < 10$ το ελάχιστο της πρώτης γραμμής
ειναι $a$ $ii)$ $8 < 10 < d$, το ελάχιστο της δεύτερης γραμμής ειναι
$8$. Για $a \in [0, 8]$ το μέγιστο εκ των δύο θα είναι το $8$, ενώ για
$a \in (8, 10)$ το μέγιστο θα είναι το $a$. Αντίστοιχα για τον παίκτη των
γραμμών θα έχουμε: $i)$ για $a \leq 8$ το μέγιστο της πρώτης στήλης θα είναι $8$
ενώ για $8 < a < 10$ το μέγιστο θα είναι το $a$, $ii)$ για $d > 10$ το μέγιστο
της δεύτερης στήλης θα είναι $d$. Συνολικά θα έχουμε:
\begin{itemize}
    \item $a \leq 8$:
        \begin{itemize}
            \item $v_1 = \max \{a, 8 \} = 8$.
            \item $v_2 = \min \{8, d \} = 8$.
        \end{itemize}
    \item $a > 8$:
        \begin{itemize}
            \item $v_1 = \max \{a, 8 \} = a$.
            \item $v_2 = \min \{8, d \} = 8$.
        \end{itemize}
\end{itemize}
Άρα για $a = 9$ (και για οποιονδήποτε αριθμό $d > 10$) δεν υπάρχει σημείο
ισορροπίας.

\subsection{Ερώτηση (ii)}

Θα βρούμε όλα τα σημεία ισορροπίας με αμιγείς και με μεικτές στρατηγικές στο
παρακάτω παίγνιο:
$$
  \begin{bmatrix}
    0, 0 & 5, 2 & 3, 4 & 6, 5 \\
    2, 6 & 3, 5 & 5, 3 & 1, 0
  \end{bmatrix}
$$

\subsubsection*{Αμιγείς στρατηγικές}
Έχουμε τα εξής σημεία ισορροπίας:
\begin{itemize}
    \item $(6, 5)$: ο παίκτης γραμμών παίζει την πρώτη γραμμή ενώ ο παίκτης
        στηλών την τελευταία στήλη. Δεδομένουε ότι ο παίκτης στηλών παίζει την
        τελευταία στήλη, ο παίκτης γραμμών δεν έχει κίνητρο να αλλάξει ($6 >
        1$), και αντίστοιχα για τον παίκτη στηλών ($5 > 0, 2, 4$).
    \item $2, 6$: ο παίκτης γραμμών παίζει την δεύτερη γραμμή ενώ ο παίκτης
        στηλών την πρώτη στήλη. Δεδομένου ότι ο παίκτης στηλών παίζει την πρώτη
        στήλη, ο παίκτης γραμμών δεν έχει κίνητρο να αλλάξει ($2 > 0$), και
        αντίστοιχα για τον παίκτη στηλών ($6 > 5, 3, 0$).
\end{itemize}

Για οποιοδήποτε άλλο προφίλ, κάποιος παίκτης έχει συμφέρον να αλλάξει
στρατητική.

\subsubsection*{Μεικτές Στρατηγικές}
Ο παίκτης γραμμών έχει μόνο δύο στρατηγικές που μπορεί να επιλέξει με πιθανότητα
$p$ και $1-p$ αντίστοιχα. Υπολογίζουμε την αναμενόμενη ωφέλεια του παίκτη στηλων
όταν ο παίκτης γραμμών παίζει με την στρατηγική $\textbf{p} = (p, 1-p)$:

$$
\begin{cases}
    u_2(\textbf{p}, e^1) = 0p + 6 (1-p) = 6 - 6p    \\
    u_2(\textbf{p}, e^2) = 2p + 5 (1-p) = 5 - 3p    \\
    u_2(\textbf{p}, e^3) = 4p + 3 (1-p) = 3 + p     \\
    u_2(\textbf{p}, e^4) = 5p + 0 (1-p) = 5p
\end{cases}
$$

Θέλουμε ο παίκτης στηλών να μεγιστοποιεί την αναμενόμενη ωφέλειά του, οπότε μας
ενδιαφέρει η άνω περιβάλλουσα του Σχήματος~\ref{fig.expected-utility}. Επίσης,
εφόσον ο παίκτης γραμμών παίζει με μεικτή στρατηγική, μας ενδιαφέρουν μόνο τα
σημεία τομής, όπου ο παίκτης στηλών είναι αδιάφορος ως προς των στρατηγικών που
παίζει. Οποιοδήποτε άλλο σημείο, θα δήλωνε αμιγής στρατηγική για τον παίκτη
στηλών και άρα μοναδικό best response για τον παίκτη γραμμών.

\begin{figure}[h]
    \centering
    \includegraphics[width=0.5\textwidth]{expected-utility.pdf}
    \caption{Αναμενόμενη ωφέλεια του παίκτη στηλών}
    \label{fig.expected-utility}
\end{figure}

Από τις ισότητες υπολογίζουμε τις υποψήφιες τιμές του $p$
$$
\begin{cases}
    u_2(\textbf{p}, e^1) = u_2(\textbf{p}, e^2) \\
    u_2(\textbf{p}, e^2) = u_2(\textbf{p}, e^3) \\
    u_2(\textbf{p}, e^3) = u_2(\textbf{p}, e^4) \\
\end{cases}
    \Leftrightarrow
\begin{cases}
    6 - 6p = 5 - 3p \\
    5 - 3p = 3 + p \\
    3 + p = 5p
\end{cases}
    \Leftrightarrow
\begin{cases}
    p = \frac{1}{3} \\
    p = \frac{1}{2} \\
    p = \frac{3}{4}
\end{cases} \,,
$$

που μας δίνουν τις υποψήφιες στρατηγικές του παίκτη γραμμών:
\begin{itemize}
    \item $p = \frac{1}{3}$: $\textbf{p} = \left(\frac{1}{3}, \frac{2}{3}\right)$
    \item $p = \frac{1}{2}$: $\textbf{p} = \left(\frac{1}{2}, \frac{1}{2}\right)$
    \item $p = \frac{3}{4}$: $\textbf{p} = \left(\frac{3}{4}, \frac{1}{4}\right)$
\end{itemize}

Για κάθε μία από τις στρατηγικές πρέπει να βρούμε την στρατηγική
$\textbf{q}$ του παίκτη στηλών που κάνει τον παίκτη γραμμών αδιάφορο. Για κάθε
υποψήφια στρατηγική, ο παίκτης στηλων χρησιμοποιεί μόνο 2 στρατηγικές, οπότε
αρκεί να βρούμε την πιθανότητα $q$ και θα παίζει την άλλη με $1-q$.

\paragraph{$p = \frac{1}{3}:$}
Έχουμε $\textbf{q} = (q, 1-q, 0, 0)$. Θα πρέπει να ισχύει
$$
    u_1(e^1, \textbf{q}) = u_1(e^2, \textbf{q}) \Leftrightarrow
    0q + 5 (1 - q) = 2q + 3 (1 - q) \Leftrightarrow
    q = \frac{1}{2} \,,
$$

και άρα έχουμε σημείο ισορροπίας το
$(\textbf{p}, \textbf{q}) = \left((\frac{1}{3}, \frac{2}{3}), \left(\frac{1}{2}, \frac{1}{2}, 0, 0\right)\right)$.

\paragraph{$p = \frac{1}{2}:$}
Έχουμε $\textbf{q} = (0, q, 1 - q, 0)$. Θα πρέπει να ισχύει
$$
    u_1(e^1, \textbf{q}) = u_1(e^2, \textbf{q}) \Leftrightarrow
    5q + 3 (1 - q) = 3q + 5 (1 - q) \Leftrightarrow
    q = \frac{1}{2} \,,
$$

και άρα έχουμε σημείο ισορροπίας το
$(\textbf{p}, \textbf{q}) = \left((\frac{1}{2}, \frac{1}{2}), \left(0, \frac{1}{2}, \frac{1}{2}, 0\right)\right)$.

\paragraph{$p = \frac{3}{4}:$}
Έχουμε $\textbf{q} = (0, 0, q, 1 - q)$. Θα πρέπει να ισχύει
$$
    u_1(e^1, \textbf{q}) = u_1(e^2, \textbf{q}) \Leftrightarrow
    3q + 6 (1 - q) = 5q + (1 - q) \Leftrightarrow
    q = \frac{5}{7} \,,
$$

και άρα έχουμε σημείο ισορροπίας το
$(\textbf{p}, \textbf{q}) = \left((\frac{3}{4}, \frac{1}{4}), \left(0, 0, \frac{5}{7}, \frac{2}{7}\right)\right)$.
